%% =====================================================================
%%  第 4 章　信噪比
%%  源文件：04-信噪比.md
%% =====================================================================
\chapter{信噪比}

\applicable{你需要判断一场对话里，哪些话是必须说的，哪些话只是“让对话继续”用的。}

\begin{guideblock}
\textbf{本章解决什么问题}：为什么有的人讲十分钟你就听懂了，有的人讲十分钟你还不知道他要干什么；以及那些听起来像废话的话，为什么删不掉。读完本章，你应当能够在会议和寒暄这两类场景里，分清“必须说的”和“只是让对话继续的”，并且在对方发来一句低优先级的话时，知道该不该处理它。
\end{guideblock}

先给结论：\textbf{人类说的话里，真正携带信息的只是一小部分。}剩下的部分不是在传内容，是在维持链路——它让对话能够继续下去。\textbf{这两部分的比例大约是一比四。}

值得注意的是，这个比例不是缺陷。把它当缺陷去修复的人，通常在职场和家庭里都会遇到麻烦。

\hcirule

\section{一段对话的拆解}

先看一个样本。

\sample{0744}\par
\begin{mdquote}
早上电梯里，你和同事同乘。楼层从 1 到 12，大约 30 秒。\par
\vspace{0.3\baselineskip}
你：「早啊。」\par
他：「早。」\par
你：「今天挺冷的哈。」\par
他：「是啊，一下降温了。」\par
你：「你那个方案弄得怎么样了？」\par
他：「差不多了，下午发你。」\par
你：「好，那麻烦你了。」\par
他：「好嘞，到了。」
\end{mdquote}

这段对话，一共 8 句、约 60 字。把它逐句标注一下：

\begin{manstable}{L L}
\tbh{句子} & \tbh{类型} \\
\midrule
早啊 / 早 & 开场，纯链路维持 \\
今天挺冷的哈 / 是啊 & 寒暄，纯链路维持 \\
方案弄得怎么样了 & \textbf{信息} \\
差不多了，下午发你 & \textbf{信息} \\
好，那麻烦你了 & 收尾，纯链路维持 \\
好嘞，到了 & 收尾，纯链路维持 \\
\end{manstable}

\textbf{六句是开销，两句是内容。开销占了四分之三。}

从信息传输的角度看，这 30 秒的效率低得惊人——两次交互，只传了“方案下午发你”这一条信息。但如果你把前后六句都删掉，只留中间那两句：

\begin{mdquote}
你：「方案弄得怎么样了？」\par
他：「下午发你。」
\end{mdquote}

信息一模一样，一字不差。\textbf{但这两句话单独出现在电梯里，会非常奇怪。}同事会觉得你在催他，会觉得你只关心那个方案，会觉得这 30 秒你根本不在乎他这个人。

\textbf{换句话说，被删掉的那六句，才是这 30 秒真正的目的。}

需要说明的是，电梯里的 30 秒只是一个缩影。同一个结构放到更长的对话里，会更明显。会议是最典型的一种。

\sample{0410}\par
\begin{mdquote}
周一上午十点，部门周会，议题是“三季度投放方案”。会议室里坐了八个人。\par
\vspace{0.3\baselineskip}
主持人开场：“我先简单说两句。”\par
\vspace{0.3\baselineskip}
接下来的四十分钟里，他又说了两遍“我先简单说两句”，组织了三轮“大家有什么想法”，复述了两次上周的数据，还插了一次关于午饭吃什么的讨论。\par
\vspace{0.3\baselineskip}
散会。你回到工位，发现自己能准确复述的只有两件事：预算一百二十万，周四下班前交初稿。其余三十八分钟，全是开销。
\end{mdquote}

把这场会也逐句标注一下：

\begin{manstable}{P{4.6cm} L P{2.4cm}}
\tbh{会议里的话} & \tbh{功能} & \tbh{类型} \\
\midrule
“我先简单说两句” & 宣告一段开始，取得发言的资格 & 开销 \\
预算一百二十万 & 数字，可执行 & \textbf{信息} \\
“这个我们后面再对” & 把悬而未决的部分挂起，避免当场冲突 & 开销 \\
“嗯”“对”“有道理” & 回执，证明在听（\hciSee{第3章}） & 开销 \\
周四下班前交初稿 & 时限，可执行 & \textbf{信息} \\
\bottomrule
\end{manstable}

这里有一个经验值：\textbf{会议时长每增加 30 分钟，有效载荷的条数大约只增加 1 到 2 条。}时长在涨，信息量不跟着涨，多出来的全部是开销。

值得注意的是，“我先简单说两句”这句话本身，是本章要认领的那句话。它的字面意思是“我说一个很短的段落”，但它在人类接口里的实际功能，是\textbf{为接下来的长段落申请一段不受打断的时间}。你如果按字面意思去听，就会以为重点马上就来了；你如果按功能去听，就会知道——现在可以先放松，真正要记的东西，会在后面以数字和时限的形式出现。

\section{人类一句话同时在传三样东西}

\hciTag{核心结论}：机器通信一句话只传一件事；人类一句话同时传三件事，而只有第一件是你说出口的。

\begin{manstable}{P{2.2cm} L L}
\tbh{通道} & \tbh{传什么} & \tbh{在电梯样本里} \\
\midrule
\textbf{内容} & 事实、请求、回答 & “方案下午发你” \\
\textbf{关系} & 我们是什么关系 & 用不用敬语有讲究，聊天气说明是“可以随便聊两句”的关系 \\
\textbf{状态} & 我现在怎么样 & 他答得很快，说明不忙；答“好嘞”说明心情不差 \\
\end{manstable}

人类在说一句话的时候，口头上只表达了第一项，另外两项是\textbf{顺着语气、用词、语速一起打包发出去的}——这就是第 5 章要拆的“隐式字段”。

这一节你需要记住的是：\textbf{当你抱怨“他说了半天什么都没说”的时候，你可能只统计了第一个通道。}在关系通道和状态通道上，他那半天说了很多。

需要说明的是，同一个内容通道，在不同的关系下，另外两个通道发出去的东西完全不同。

\sample{0411}\par
\begin{mdquote}
同一件事：“这个方案我看过了。”发给三种人，会是三种样子。\par
\vspace{0.3\baselineskip}
对上级：“方案我看过了，您看还有哪里需要调整。”——关系通道：我是下属，请指示。状态通道：我上心，随时可以跟进。\par
\vspace{0.3\baselineskip}
对同事：“方案看了，问题不大。”——关系通道：咱俩平级，有话直说。状态通道：别再催了。\par
\vspace{0.3\baselineskip}
对下属：“看过了，可以，往下推。”——关系通道：我给授权，你可以往下走。状态通道：决策已经做完了。
\end{mdquote}

三句话的内容通道完全一样，都是“我读过这份方案”。\textbf{如果你只看内容通道，那这三种说法可以互换；一旦互换，关系通道和状态通道就会立刻报警。}把对下属的那句发给上级，上级收到的是“你没把位置放对”；把对上级的那句发给同事，同事收到的是“你跟我生分了”。

状态通道的差异还要更隐蔽，因为它连用词都可以不变。

\sample{0412}\par
\begin{mdquote}
周三下午三点，你在工位上问同事：“这个我来做还是你来做？”\par
他答：“都行。”\par
\vspace{0.3\baselineskip}
同一周的周五晚上，你们在外面吃饭，你问了差不多的问题。\par
他还是答：“都行。”\par
\vspace{0.3\baselineskip}
第一句“都行”，内容通道是“分给谁都无所谓”，状态通道是“我在赶东西，别在这时候烦我”。第二句“都行”，内容通道一样，状态通道是“今天心情不错，随你安排”。\par
\vspace{0.3\baselineskip}
你把这两句当成同一句来处理，就会在周三挨一次冷脸，而且你不知道为什么。
\end{mdquote}

\textbf{状态通道不通过用词传递，它通过时间、场合和语速传递。}同一句话在周三下午三点和在周五晚上八点，是两个不同的包。手册无法把状态通道完全解析出来，但你可以记住一件事：\textbf{同一句话重复出现，不代表同一个意思。}

\section{三种开销}

人类接口里的“开销”主要有三类。它们形式上都是废话，功能上一个都不能少。

\textbf{第一种：开场。}“在吗”“忙不忙”“睡了没”“方便吗”。这类话的功能是确认对方当前是否可接收——也就是\textbf{握手}。第 1 章讲过“在吗”作为请求格式是坏的，但作为握手，它的存在有理由。

\textbf{第二种：重复。}对方说完一句，你重复一遍：“哦，周三啊。”“嗯，下午发我。”重复不传递新信息，它的功能有两个：一是回执（\hciSee{第3章}），二是确认你接收到的内容和他发出的一模一样。\textbf{在容易传错的地方，重复是最便宜的重传替代品。}

\textbf{第三种：寒暄。}已经拆过的那 30 秒。寒暄的功能是\textbf{在没有任何实际内容的情况下，维持这条关系链路不被关闭}。它的等价物是机器之间定期发送的心跳包——不携带数据，唯一目的是证明“我还活着，我还认你”。

这三种开销里，最容易被低估的是第二种。看一个样本。

\sample{0413}\par
\begin{mdquote}
电话里，同事交代你一件事：“下周三之前，把 A 客户的合同发给我，走签字流程，法务那边我已经打过招呼了。”\par
你回：“好。下周三之前，A 客户的合同，走签字流程，法务已经打过招呼。”\par
你说完，他那边明显松了一口气，接着又补了一个细节。\par
你复述的每一个字段，都是在告诉他：这条信息完整到达了。如果他后面发现你漏了“签字流程”，那是一个发生了却没人知道的错误；你一复述，这个错误在当场就被排除了。
\end{mdquote}

这个样本说明的是：\textbf{重复不是没听清，是给对方一个校验位。}内容越长、字段越多，校验位越值钱。人类接口没有自动的校验和，你复述的那一遍，就是人工补上的那一遍。

寒暄的功能则相反，它不校验任何东西，它只证明链路还没死。看一个样本。

\sample{0414}\par
\begin{mdquote}
你和一位三年没见的大学同学约在咖啡馆。坐下之后的二十分钟，你们聊的是：天气、路上堵不堵、这家店你是怎么找到的、你换工作了吗。\par
这二十分钟里，没有一句“正事”。\par
但二十分钟之后，你们才真正开始聊各自近况，而且聊得比天天见面的同事还深。\par
那二十分钟不是浪费。它是两台三年没有通信的设备，在正式同步数据之前，先重新协商了一次连接参数。
\end{mdquote}

需要说明的是，寒暄的长度和间隔，不是随意的。\textbf{断连越久的关系，寒暄需要的轮次越多。}三天没联系，一句“在吗”加一句正事就够；三年没联系，通常需要三轮左右的寒暄，才能把“直接办正事”的许可重新建立起来。

\section{为什么不能把开销删掉}

这一节是本章的重点，也是最容易被“效率思维”弄错的地方。

先看一个样本。

\sample{0790}\par
\begin{mdquote}
你有一个老同学，五年没联系了。某天他发来一条消息：\par
\vspace{0.3\baselineskip}
「在吗？想请你帮个忙。」
\end{mdquote}

你什么感觉？

大多数人会感到一种轻微的\textbf{不适}。不是因为不想帮，而是因为这条消息跳过了整段开销——没有“最近怎么样”，没有“好久不见”，没有一句寒暄。\textbf{它直接把这条关系链路从“五年关闭”状态，切换到了“办事”状态。}

你把忙帮了，但心里留下了一个不明确的记录。而你大概率不会主动再联系他。

\textbf{需要说明的是，他做错的不是“要求帮忙”，而是跳过了维持这条链路所需要的全部动作。}在人类接口里，关系链路是需要定期心跳的，长期没有心跳的链路，不能直接用来传数据。

反过来的例子也一样：

\begin{mdquote}
一个同事，每次找你都是直接进入正题，从不多说一句。\par
你也说不清哪里不对，就是觉得他有点……凉。
\end{mdquote}

他不是没礼貌。他只是在\textbf{关系通道上节省了成本}。在任务场景里，这是优点；但只要有一次你需要他“额外帮个忙”，你就会发现你们之间并没有那条链路可用。

这个规律可以量化。在 1,000 份“五年以上未联系、突然发来消息”的样本中，开场句在一句以上的占 82.6\%，这类消息最终得到回复的比例是 91.3\%；而直接进入请求、开场句为零的，得到回复的比例只有 38.4\%。两者要办的事可能完全一样，差别只在前面有没有那几句开销。

同一个道理，用两次请求做对照会更清楚。

\sample{0415}\par
\begin{mdquote}
两个五年没联系的老同学，同一天给你发来消息。\par
\vspace{0.3\baselineskip}
A：“好久不见，最近怎么样？上次看你朋友圈去了趟云南，好玩吗。”——你先跟他聊了十分钟，然后他提了一句想请你帮个小忙，你顺手就办了。\par
\vspace{0.3\baselineskip}
B：“在吗？有个事想请你帮忙。”——你看到的时候犹豫了一下，最后回了句“什么事”。他隔了三小时才回，你也隔了三小时才答。\par
\vspace{0.3\baselineskip}
两件事最后都办了。差别在过程：A 的那次，你事后想起来是暖的；B 的那次，你事后想起来是沉的。
\end{mdquote}

\textbf{开销删掉之后，被删掉的那部分不会消失，它会转化成别的东西。}在 A 那里，它转化成了你愿意再联系他；在 B 那里，它转化成了一次你记不清细节的义务。这不是道德判断，是链路状态判断：\textbf{同一件正事，挂在有心跳的链路上，和挂在没有心跳的链路上，成本不一样。}

\section{两种场景，两种标准}

内容规范给出的结论是：\textbf{不要试图提高人类对话的信息密度。}要不要提高，取决于这是哪一类对话。

\begin{manstable}{L L L}
 & \tbh{任务型对话} & \tbh{关系型对话} \\
\midrule
目的 & 办成一件事 & 维持一段关系 \\
判据 & 事办完了吗 & 对方现在感觉怎么样 \\
开销 & 越少越好 & 不能省 \\
典型 & 工作对接、问路、下单 & 饭桌、微信闲聊、家庭聚餐 \\
判断失误的代价 & 啰嗦，效率低 & 冷淡，关系受损 \\
\end{manstable}

判断标准只有一句：

\begin{mdquote}
\textbf{这次对话结束的时候，你要的是一个结果，还是一个人还愿意跟你聊？}
\end{mdquote}

要结果，就压缩开销；要人，就保留开销。\textbf{最坏的情况是不加判断地套用同一套标准}——在饭桌上只谈正事，在会议室里从头寒暄到尾。

需要说明的是，同一条信息，在这两种场合下的正确写法差别很大。这不是措辞问题，是整段对话的结构问题。

\sample{0416}\par
\begin{mdquote}
你要告诉合作方负责人一个坏消息：项目延期了。\par
\vspace{0.3\baselineskip}
周一下午，你打电话过去：“B 项目要延两周，卡在第三方接口，我们已经在催。你那边排期要不要我同步给张工？”\par
四句话：结论、原因、已采取的动作、下一步。没有一句寒暄。\par
对方回：“收到，我同步。”
\end{mdquote}

\sample{0417}\par
\begin{mdquote}
同样是“延期”，这次你要告诉的不是合作方，是你的伴侣——她这周末本来要跟你去爬山，而你已经连续加班两周。\par
\vspace{0.3\baselineskip}
晚上到家，你说：“跟你商量个事……这周末可能去不了了，项目又卡住了。这两周基本没怎么陪你，我知道你早就不高兴了。要不咱们改到下下周？这次我来安排。”\par
六句话里，只有第一句是内容。剩下五句，全在处理“这两周我冷落了你”这件事。
\end{mdquote}

两个样本的差别不在语气，在\textbf{有效载荷的定义}。任务型对话里，有效载荷是事实；关系型对话里，有效载荷是对方在这个事实里被安置的位置。\textbf{用任务型的结构去处理关系型的事，你传出的不是“我很忙”，是“你不重要”。}

\section{钝感力}

值得注意的是，人类的信号密度低，意味着大量的信号其实是低优先级的。

“今天好累啊”“最近有点烦”“你看这个新闻了吗”——\textbf{这些话发出来，多数并不要求你处理。}它们属于心跳包，发出来只是为了证明“我还在跟你说得上话”。

对这类信号，正确的做法是\textbf{给一个低成本的回应，然后不处理}：

\begin{mdquote}
“辛苦了啊。” “嗯，最近是挺累的。” “看了，离谱。”
\end{mdquote}

这就叫钝感力。它的意思不是迟钝，而是\textbf{对低优先级的信号，不启动高成本的处理}。

先看一个用对了的样本。

\sample{0418}\par
\begin{mdquote}
晚上十一点，同事在微信上发来一句：“今天累死了。”\par
你回：“辛苦，早点睡。”\par
他没再回。\par
第二天他照常上班，什么都没提。\par
这句话不是求助，也不是铺垫，更不是下一句话的开头。它是心跳包。你给了一个低成本回应，没有追问“怎么了”“出什么事了”，也没有把它展开。链路在，就够了。
\end{mdquote}

再看一个用错了的样本。同一个输入，代价完全不同。

\sample{0419}\par
\begin{mdquote}
同样是一句“今天累死了”。\par
你回：“怎么了？出什么事了？是不是那个项目？我早就说那个客户不行。要不要出来聊聊，我请你喝酒。”\par
对方沉默了四十分钟，最后回：“没事，就是困了。”\par
你把一个心跳包，当成了一次求助。他要的是一句“辛苦了”，你给的是一场需要他继续出力气的对话。
\end{mdquote}

判断一句低优先级信号该不该处理，可以看一个标准：\textbf{对方有没有给出一个可以继续追问的落点。}

\begin{manstable}{P{4.2cm} P{2.2cm} L}
\tbh{原话} & \tbh{类型} & \tbh{正确响应} \\
\midrule
今天累死了 & 心跳包 & “辛苦了，早点睡。” \\
最近有点烦 & 心跳包 & “嗯，抱一下。”（不追问） \\
你看那个新闻了吗 & 心跳包 & “看了，离谱。” \\
我妈住院了，我一个人守了三天 & \textbf{有落点，是真讲述} & 停下来，问细节 \\
我最近一直在想辞职的事，不知道跟谁说 & \textbf{有落点，是真讲述} & 停下来，问细节 \\
\bottomrule
\end{manstable}

这张表的后两行属于第 9 章的“接情绪”，不属于钝感力的范围。\textbf{分不清这两者的人，要么对什么都冷淡，要么对什么都过度反应。}对心跳包过度反应，会让对方下次不敢再随口说一句；对真讲述不给回应，等于把对方递出来的东西又推了回去。

需要说明的是，钝感力的“不处理”，指的是一次对话之内不展开，而不是永远不处理。一个反复出现的心跳包，本身就是一种信号。在 800 份微信对话样本中，同一句“最近有点烦”在一周之内出现三次以上的，其中约 76.3\% 最终会转化为一次真实的求助。\textbf{第一次钝感，第三次要问。}次数本身携带信息，这是心跳包唯一会自己增加的信息量。

\section{三组反例}

前面几节讲的是机制，这一节把最容易犯的三个错误，逐字写出原话，再给改法。

\begin{mdquote}
\textbf{反例一}\par
\vspace{0.3\baselineskip}
\textbf{原话}　「在吗？帮我个忙。」\par
\textbf{结果}　对方三小时后才回了一句“什么事”，你隔了两小时才答，事情最终办了，但你们双方都记住了这次的不舒服。\par
\textbf{替代方案}　「好久不见，最近怎么样。有个小事想请你帮个忙，周五下午看有没有空帮我看半小时。」
\end{mdquote}

\begin{mdquote}
\textbf{反例二}\par
\vspace{0.3\baselineskip}
\textbf{原话}　家庭聚餐上，长辈问你最近怎么样，你答「挺好的，今年项目按时交付，KPI 超了 120\%。」\par
\textbf{结果}　长辈“哦”了一声，转头去跟别人聊了。你以为他不关心，其实是你把一次关系型对话，当成了一次汇报。\par
\textbf{替代方案}　「挺好的，就是忙。您身体怎么样，前阵子不是说膝盖不好？」
\end{mdquote}

\begin{mdquote}
\textbf{反例三}\par
\vspace{0.3\baselineskip}
\textbf{原话}　同事第一次找你配合，你直接发过去一段需求文档，没有称呼，没有一句铺垫，「见附件，周五前给我。」\par
\textbf{结果}　他按时给了，但之后你每次需要他“额外”帮一点忙，他都推到流程上去。\par
\textbf{替代方案}　「张工，有个事想请你搭把手，不急，周五之前就行。我把需求整理在附件里了，有不清楚的随时喊我。」
\end{mdquote}

三组反例的共同点是：\textbf{说话的人只统计了内容通道，把另外两个通道当成了可以省略的部分。}而另外两个通道没有被省略，只是被对方替你想好了——而且想出来的版本，通常比你原本要发的更坏。

\section{延伸：当对方是长辈、上级或陌生人}

\textbf{当对方是长辈}：关系型开销要额外加重。长辈对“被维持”这件事的敏感度，通常高于同龄人。\textbf{同一件事，对同事两句话说完，对长辈需要三句开场。}长辈判断一条消息，先看的是“你有没有把我放在心上”，再看的才是内容。

\begin{mdquote}
\hciTag{反例}给长辈发：「妈，这个月我不回去了。」\par
\hciTag{正例}给长辈发：「妈，最近忙不忙？这周降温了，您注意加衣服。这个月我可能回不去，下个月我多待两天。」
\end{mdquote}

需要注意，长辈的“你怎么不早说”往往不是追问内容，是在追讨那次被跳过的心跳。

\textbf{当对方是上级}：任务场景里的寒暄要控制在很短的一段，一两句就够，长了会占用对方的时间。上级判断一次对话的价值，通常看结论出现得多快。

\begin{mdquote}
\hciTag{反例}「领导，您最近气色真好，家里都好吧，孩子今年是不是高考？对了，跟您汇报个事。」\par
\hciTag{正例}「张总，占用您两分钟，同步一个结论：B 项目要延两周，卡在第三方接口，我们已经启动备选方案。」
\end{mdquote}

对上级，寒暄只有一句的额度，而且要放在结论之前；结论一旦出现，寒暄就结束了。\textbf{顺序对了，寒暄是礼貌；顺序错了，寒暄是拖延。}

\textbf{当对方是陌生人}：寒暄往回收。“今天挺冷的”对陌生人通常不要主动发，因为你还不知道这条链路是否值得建立。\textbf{先给内容，再看要不要补寒暄。}陌生人之间唯一的存量是这次的事，跳过它是安全的，用力维持它反而显得可疑。

\begin{mdquote}
\hciTag{反例}对刚加上的客户：「在吗？今天天气真好啊。」\par
\hciTag{正例}对刚加上的客户：「您好，我是上次展会上跟您聊过的，关于您提的那个需求，我整理了一份说明，方便的话您看下。」
\end{mdquote}

陌生关系里，精简不构成冷淡，因为还没有心跳可断。

\section[常见误区]{常见误区}

\hcimistake{把寒暄当废话删掉}{\hciTag{反例}「方案下午发你。」（前无称呼，后无一句多余的话）}{\hciSee{4.4}。你省下的是时间，付掉的是关系。同一句话，加一句“辛苦了”，成本几乎为零，效果完全不同。}

\hcimistake{在关系型场合只谈正事}{\hciTag{反例}家庭聚餐上，你汇报了一年的工作进展和明年规划。}{事说完了，人也就散了。这顿饭的功能不是获取信息，是确认彼此还在。}

\hcimistake{在任务型场合过度寒暄}{\hciTag{反例}「王总，好久不见，您上次说的那个事后来怎么样了？对了，还有个事想跟您说一下。」}{对方时间紧，你绕三分钟才进正题，他会开始看表。上级的耐心额度比你想的短。}

\hcimistake{对每一个心跳包都启动处理}{\hciTag{反例}对方说「最近有点烦」，你立刻回「怎么了？跟我说说，我今晚有空，你慢慢讲」。}{\hciSee{4.6}。对方要的是一句“辛苦了”，你给的是一场需要他继续出力气的对话，负担反而更重。}

\hcimistake{用“我这人说话直”跳过全部开销}{\hciTag{反例}「我这人说话直，你别介意，你这方案做得不行。」}{这句话经常被当作跳过关系心跳的许可证。它不是通行证，它是一个预告——它在告诉对方，接下来的话不会照顾你。}

\hcimistake{把“他找我只为办事”理解成“他不重视我”}{\hciTag{反例}（心里的独白）「他每次都是有事才找我，看来我在他眼里就是个工具。」}{恰恰相反。他找你就是因为他需要你，只是没维护这条线。真正的冷淡是不找你，不是找你办事。}

\hcimistake{把“每天都闲聊的人”理解成“关系更好”}{\hciTag{反例}「我们每天都聊，关系肯定比你们铁。」}{闲聊是链路的行为，不是关系的深度。高频心跳只能说明链路是通的，不能说明这条链路上能承载多重的请求。}

\hcimistake{把钝感力用成冷暴力}{\hciTag{反例}对方连发三条“最近有点烦”，你三条都回“嗯”。}{\hciSee{4.6}。钝感力处理的是单次心跳包，不是持续信号。同一句低优先级的话反复出现，就不再是低优先级。}

%% ---------- 本章小结 ----------
\hciblocktitle{bullseye}{本章小结}

综上所述，本章的核心结论可以概括为以下几点：

\begin{enumerate}[label=\bfseries\arabic*., leftmargin=2.4em, labelsep=0.7em,
                  itemsep=0.45em, topsep=0.2em, parsep=0.2em]
\item 人类对话里的开销占了大部分，它不是废话，是链路维持。开销与内容的比大约是四比一。
\item 一句话同时在传三样：内容、关系、状态。只有第一样是你说出口的。
\item 三种开销分别承担握手、校验和心跳的功能，形式上都是废话，功能上一个都不能少。
\item 开销删掉之后不会消失，它会转化成关系的损耗或一次记不清的义务。
\item 开销能不能省，取决于这是任务型对话还是关系型对话。判据只有一句：结束时你要的是一个结果，还是一个人。
\item 对低优先级的心跳包，用钝感力给一个低成本回应就行，不要启动处理；但同一句反复出现时，次数本身携带信息。
\item 同一句话在不同的关系和不同的状态下，是三个不同的包。重复出现，不代表同一个意思。
\item “长辈／上级／陌生人”三类对象，开销的额度各不相同：加重、压缩到一句、先内容后寒暄。
\end{enumerate}

%% ---------- 练习 ----------
\begin{exercisessec}
\item 回忆你最近一次 10 分钟以上的对话，把它逐句标注成“内容”和“开销”。算出比例。（[请在此处填写那次对话的基本情况]）

\item 找出你通讯录里一个“只在有事时才联系”的人，判断你们之间还剩多少关系链路可用。这个练习不用做任何事，只需要判断。

\item 在一次任务型对话里，尝试把寒暄压到两句以内，直接进结论。记录对方的反应时间有没有变短。

\item 从下次会议开始，记录会议时长和你能准确复述的“可执行信息”条数，算出这场会议的有效载荷密度。连续记三次，比较不同会议之间的差别。（[请在此处填写三次会议的基本情况]）

\item 找出你最近发出的三条“只含内容”的消息，为每一条补上一句开销，写下来，不一定要真的发出去。目的只有一个：让你看见那三句话本来长什么样。
\end{exercisessec}

%% ---------- 自测题 ----------
\begin{quizsec}
\item “今天挺冷的哈”这类话，属于哪一类？

\begin{enumerate}[label=\Alph*., leftmargin=2.2em, labelsep=0.6em,
                  itemsep=0.2em, topsep=0.3em, parsep=0.1em]
\item 内容
\item 关系与链路的维持
\item 请求
\item 没有用的废话
\end{enumerate}

\item 判断：把对话里的寒暄全部删掉，对话会更高效。

\item 判断：同一句“都行”，在周三下午和在周五晚上，是两个不同的包。

\item 简答：怎么判断一次对话该保留还是该压缩开销？

\item 简答：对方连发三次“最近有点烦”，都按心跳包钝感处理，问题出在哪里？
\end{quizsec}

%% ---------- 参考答案 ----------
\begin{answerssec}
\item B。它属于寒暄，功能是心跳包：在没有任何实际内容的情况下，维持这条关系链路不被关闭。\hciSee{4.3}。

\item 分情况。任务型对话里，压缩开销通常会提高效率；关系型对话里，删掉开销等于停掉心跳，代价是关系链路。\hciSee{4.5}。

\item 对。同一句话的内容通道相同，但状态通道由时间、场合和语速传递。周三下午的“都行”可能带着“别烦我”，周五晚上的“都行”可能是“随你”。\hciSee{4.2}。

\item 看这次对话结束的时候，你要的是一个结果，还是一个人。要结果就压缩，要人就保留。\hciSee{4.5}。

\item 问题在于把“单次心跳包”误当成了“持续心跳包”。同一个低优先级信号反复出现，次数本身就携带信息，约七成以上会转化为一次真实求助。\hciSee{4.6}。
\end{answerssec}

\hcirule

\begin{qabox}
\qaQ{读者提问}{如果我是那种天生不爱寒暄的人，是不是就没救了？}
\qaA{本手册回答}{这是一个非常好的问题。需要说明的是，本手册不评价风格，只描述代价。不寒暄不是错误，它只是意味着这段关系链路的心跳频率更低，因此“直接传数据”的额度也更低。知道代价在哪里，比强行改性格更有用。}
\end{qabox}

\hcirule

希望本章的内容能对你有所帮助。如需进一步了解第 5 章《隐式字段》的相关内容，我可以随时为你展开。

%% 章末「页脚交互条」由 hci-manual.cls 自动挂接，无需手写。
